\originalsection{电阻为何从侧面接收场能}
设一段长直圆柱导线半径 $a$、长度 $\ell$, 导电率 $\sigma$, 直流电流 $I$ 沿 $+z$. 在远离端点的均匀电流近似中, $E_z=I/(\sigma\pi a^2)$; 导线侧面切向电场连续, 表面磁场 $B_\phi=\mu_0I/(2\pi a)$. 所以侧面的 $\vct S=\ee\times\bb/\mu_0$ 指向径向内侧, 因 $\hat{\vct z}\times\hat{\vct\phi}=-\hat{\vct r}$.

单位时间从侧面进入的总能量为
$2\pi a\ell\,E_zB_\phi/\mu_0=I E_z\ell=I^2\ell/(\sigma\pi a^2)$.
体内焦耳功率积分也给
$\int\jj\cdot\ee\dd^3x=\sigma E_z^2\pi a^2\ell=I^2R_\mathrm{wire}$,
$R_\mathrm{wire}=\ell/(\sigma\pi a^2)$. 静态场并不意味着没有能流; 这里 $\partial_tu=0$, 入流立即转成热. 维持电流的外电路负责建立上述场, 端部及返回导线附近需用完整电路几何处理. 这个局部计算解释本科作业9只给教材题号的相关概念, 不冒充缺失的教材题面.
